

\illustration{0.28}{img/cowsay.png}{}%

Karl is an aspiring C programmer, and is excited by the risks and rewards of
low-level manual memory management.
In the program he currently develops, he stores a string containing $N$ non-zero
bytes into a buffer named ``\texttt{buf}''.
By mistake he accidentally made the buffer $2 N$ bytes in size. The last $N$
bytes of the buffer consists of only zero-bytes.

\medskip

\noindent
Now Karl needs to know the value $N$, the size of the string, in a separate part
of the program. Traditionally you would recover the length of a string using the
\texttt{strlen}-function, which reports the position of the first zero-byte in
the provided buffer using a linear scan.
However, Karl finds that this is much too slow, and that it defeats the advantage
of using C in the first place.
Can you help Karl efficiently recover $N$ without crashing his program?

\begin{center}
\includegraphics[width=.6\textwidth]{img/buffer.pdf}

\textsl{The contents of the buffer in sample interaction 3 are shown here.}
\end{center}

